\section{Methods}
\label{sec:methods}

\subsection{The \jaisp{} foundation model}
\label{sec:methods:foundation}

% --- architecture ---
\jaisp{} encodes each band with a per-band convolutional stem, resamples the
per-instrument streams to a common fused pixel scale of $0.4\arcsecpx$, and
fuses them into a shared bottleneck representation. \todo{Architecture detail:
per-band stems, mixed-resolution stream depths from $\log_2$(scale ratio),
ConvNeXt/transformer blocks, FiLM band conditioning, concat fusion (v9/v10).}
The per-band stem normalizes its input by the per-pixel noise,
$x = \mathrm{image}/\mathrm{rms}$, and clamps the result to a fixed
signal-to-noise range before feature extraction; this choice is revisited in
\S\ref{sec:results} as the origin of bright-star detection suppression.

% --- training ---
The model is trained by masked reconstruction \todo{(MAE objective, masking
fraction, cross-instrument masking, Charbonnier$+$core-L2 loss for v10), epochs,
optimizer, hardware}. All downstream heads operate on \emph{frozen} features.

% --- versions / equivalence note ---
\todo{Briefly: several architecture variants (v8 fine-scale, v9 concat fusion,
v10 loss-shaped) were trained; in-distribution probes and downstream metrics
find them statistically equivalent (see \S\ref{sec:discussion}). Report the
production checkpoint used here (v10 warm-start).}

\subsection{Detection head}
\label{sec:methods:detection}

% --- centernet ---
We detect sources with a center-prediction (CenterNet-style) head on the frozen
fused features, producing a normalized centroid and score per source.
\todo{Head architecture, self-training/pseudo-label procedure, operating point
(conf threshold), NMS.} Detections are gated on the \textit{Euclid} VIS signal
and serve as seeds for astrometry.

\subsection{Multi-band joint-centroid astrometry}
\label{sec:methods:astrometry}

% --- the instrument ---
Given a detection seed, we (i) refine the centroid classically on
\textit{Euclid} VIS, (ii) solve for a single \emph{joint canonical} sky
position by simultaneously fitting the source across all bands, weighting each
band by its noise and accounting for its PSF, and (iii) measure each band's
classical centroid and record its offset from the joint canonical position.
The joint canonical position is the cross-instrument estimate; the per-band
offset is the astrometric residual we report.

% --- why this beats single-band ---
Because the fit is dominated by the band with the sharpest, highest-SNR
localization (typically \textit{Euclid} VIS), the joint position transfers that
localization to bands in which the source is faint or unresolved
(\S\ref{sec:results}). Per-source uncertainties follow from the $\chi^2$
curvature of the joint fit and a per-band $\mathrm{FWHM}/(2.355\,\mathrm{SNR})$
centroid-noise model.

% --- foundation-agnostic note ---
We emphasize that step (i)--(iii) are classical given the detection seeds; the
foundation's role in astrometry is to supply the cross-instrument detections.
A learned per-object position head was also trained and is discussed as an
ablation in \S\ref{sec:discussion}.
